% 问题重述与分析
\section{问题重述与分析}

\subsection{问题背景}
无线电频谱的有序使用关系到通信、导航和广播等业务的正常运行。传统干扰排查通常经历初步监测、区域排查和人工近场定位，最后一环耗时且可能使人员暴露于危险环境。搭载测向机、光学定位装置和激光清除装置的四足机器人，能够把这一过程改写为“感知—决策—移动—清除”的自主闭环；但示向误差、未知接收半径、未知频道以及定向源盲区，使系统同时面临几何不确定性和在线路径决策问题\cite{wang2019}。

\subsection{问题重述}
\textbf{问题一：}由若干检测点坐标和同一干扰源的示向度，在 $\pm1^\circ$ 误差下构造交会定位区域，求其直径，并判断以一组直径端点为直径的圆能否覆盖该区域。

\textbf{问题二：}已知某全向干扰源在第一检测点的示向度，设计第二检测点，使再次接收具有确定性保证，并在定位精度与移动检测耗时之间取得权衡，给出候选区域及推荐站址。

\textbf{问题三：}目标圆域内有 $10$--$16$ 个频道互异的全向干扰源，其数量、位置、频道和接收半径均未知；机器狗从原点、频道 1 出发，自主完成发现、定位、清除与不存在认证，并最小化虚拟总时间。

\textbf{问题四：}在问题三中加入数量和朝向均未知的定向干扰源。此时一次无信号还可能由定向盲区造成，需重新设计发现证书与清除策略，并完成模拟器测试。

\subsection{问题分析与总体思路}
四问共享“\textbf{不确定集合—确定性证书—代价最小化}”主线。问题一把每个示向度写成闭凸测向锥，以半平面交获得定位集合，再用凸包与旋转卡壳求直径，并借最小覆盖圆和 Jung 定理回答覆盖性。问题二先由首次正检测构造完整可能源集合，再求保证接收的第二站候选区；以误差盒经 Jacobian 传播所得的最坏位置误差和局部耗时组成 Pareto 双目标，第二次读数返回后仍调用问题一的精确集合与最小覆盖圆作清除判定。

问题三把“存在源的定位清除”和“无源频道的完备认证”同时纳入在线路径：扫描站承担发现与阴性证据积累，多站示向度顺路交会，未达清除证书的频道才补充观测。问题四的关键不再是普通距离覆盖，而是使任意可能源附近的扫描点从各个半平面方向包围该源；只有这种方向完备性成立，连续无信号才足以排除未知朝向的定向源。最终以公开接口的动作日志复核清除闭合性、时间账和几何证书。
